\documentclass[11pt,a4paper]{article}
\usepackage{curso}
\begin{document}
\portada{Lección}{4.4}{Diagnósticos de muestreo}

\noindent\textbf{Unidad 4.}\quad\textbf{Objetivos:} calentamiento, $\tau_{\text{int}}$ y tamaño efectivo de muestra; desaceleración crítica;
recocido simulado y tempering paralelo.
\textbf{Prerrequisitos:} 2.6 (ergodicidad), 4.2--4.3.

\section{Calentamiento}
Las primeras muestras recuerdan el estado inicial; se descartan. Diagnóstico práctico: varias cadenas desde inicios distintos (todo arriba, todo
abajo, aleatorio) cuyas trayectorias de $E/N$ se superponen al terminar el calentamiento (estadístico $\hat R$ de Gelman--Rubin en la práctica).

\section{Autocorrelación y tamaño efectivo}
Con la autocorrelación normalizada $\rho_k$ de un observable,
\begin{equation}
\tau_{\text{int}}=1+2\sum_{k\ge1}\rho_k,\qquad \mathrm{ESS}=\frac{n}{\tau_{\text{int}}},\qquad \sigma_{\bar x}\approx\sqrt{\frac{\mathrm{Var}(x)\,\tau_{\text{int}}}{n}},
\end{equation}
truncando la suma con una ventana autoconsistente ($k\ge5\tau_{\text{int}}$).
\begin{ejemplo}
Ising $16\times16$, Gibbs cromático, $n=38\,000$ barridos tras 2000 de calentamiento, observable $m$ \evS:
\begin{center}\begin{tabular}{lccc}\toprule
$T$ & 3,5 & 2,8 & 2,27\\\midrule
$\tau_{\text{int}}$ (barridos) & 8,5 & 33 & $\approx2200$\\
ESS & 4481 & 1163 & 17\\\bottomrule\end{tabular}\end{center}
\end{ejemplo}
\begin{intuicion}
Desaceleración crítica: cerca de $T_c$ los dominios crecen hasta el tamaño del sistema y las actualizaciones locales tardan $\sim L^z$ barridos
($z\approx2{,}17$ para Gibbs/Metropolis en 2D) en reorganizarlos. Algoritmos de clúster (Swendsen--Wang, Wolff) reducen $z$.
\end{intuicion}

\section{Recocido simulado}
Para minimizar $E$, se muestrea con $T(t)$ decreciente: caliente para explorar, frío para concentrarse en mínimos. En un vidrio de espín $\pm J$
$16\times16$ con rampa lineal $T:3\to0{,}05$ \evS: $E/N=-1{,}25$ (10 barridos), $-1{,}33$ (100), $-1{,}36$ (1000), $-1{,}37$ ($10^4$). Mejor con
enfriamientos más lentos, con rendimientos decrecientes; la longitud termodinámica (3.8) sugiere ir despacio cerca de $T_c$.

\section{Tempering paralelo}
$K$ réplicas a temperaturas $T_1<\dots<T_K$; tras barridos de Gibbs se propone intercambiar réplicas vecinas con
\begin{equation}
A_{ij}=\min\big(1,\ e^{(\beta_i-\beta_j)(E_i-E_j)}\big),
\end{equation}
que cumple balance detallado en el espacio producto $\prod_k e^{-\beta_kE(s_k)}$. Ejemplo: $\beta_i=0{,}5$, $\beta_j=0{,}6$, $E_i-E_j=20$:
$A=e^{-2}=0{,}135$. Escalera geométrica $T_k=T_{\min}(T_{\max}/T_{\min})^{k/(K-1)}$ como punto de partida.

\begin{aplicacion}
Energía de un TSU por muestra independiente $=E_{\text{barrido}}\times\tau_{\text{int}}$. Laboratorio L6: calentamiento, $\tau_{\text{int}}$, ESS, recocido y
tempering paralelo en vidrios de espín.
\end{aplicacion}

\begin{error}
\begin{itemize}[nosep]
\item Medir $\tau_{\text{int}}$ en un observable rápido. \quad Corridas cortas que subestiman $\tau_{\text{int}}$.
\item Enfriar demasiado rápido. \quad Temperaturas muy separadas en el tempering: no hay intercambios.
\end{itemize}
\end{error}

\begin{thebibliography}{9}\small
\bibitem[13]{katz} H. G. Katzgraber, Introduction to Monte Carlo methods, 2009.
\bibitem[3]{lpw} D. A. Levin, Y. Peres y E. L. Wilmer, \emph{Markov Chains and Mixing Times}, AMS, 2009.
\bibitem[19]{aadit} N. A. Aadit et al., Massively parallel probabilistic computing with sparse Ising machines, 2022.
\end{thebibliography}
\end{document}
